\section{Normalization and prior work}
\label{sec:normalization}

\subsection{Words, graphs, and reverse complement}

Let $\A$ be a finite alphabet of even cardinality $q$, equipped with a
fixed-point-free involution $a\mapsto\bar a$.  The binary specialization is
$\A=\{0,1\}$ with $\bar0=1$ and $\bar1=0$.  For a word
$w=w_0\cdots w_{k-1}$, define
\[
  \RC(w)_i=\overline{w_{k-1-i}}.
\]
The directed vertex de Bruijn graph $B(q,k)$ has vertex set $\A^k$ and
an arc $u\to v$ exactly when $u_1\cdots u_{k-1}=v_0\cdots v_{k-2}$.
Loops are retained.  Reverse complement is a directed anti-automorphism:
\[
  u\to v \quad\Longrightarrow\quad \RC(v)\to\RC(u).
\]

A set $M\subseteq\A^k$ is a \emph{decycling set} when $B(q,k)\setminus M$
is a directed acyclic graph (DAG), and is an \emph{ordinary minimum
decycling set} when its cardinality is minimum among all decycling sets.
We abbreviate the latter to MDS.  Reverse-complement invariance means
$\RC(M)=M$ as a literal set of labelled words.

The line-graph identity $B(q,k)=L(B(q,k-1))$ lets us regard each
$w=w_0\cdots w_{k-1}$ as the labelled lower-state arc
\[
  t(w)=w_0\cdots w_{k-2}\longrightarrow
  h(w)=w_1\cdots w_{k-1}.
\]
Deleting $M$ as vertices of $B(q,k)$ is equivalent to deleting those
labelled arcs from $B(q,k-1)$; acyclicity is equivalent in the two
representations.

\subsection{Pure cycling registers and the ordinary minimum}

Let $R(w_0\cdots w_{k-1})=w_1\cdots w_{k-1}w_0$.  A pure cycling register
(PCR) is an orbit under $R$, using distinct rotations only.  Its cardinality
is the minimal period of its words.  PCRs partition $\A^k$, and each is a
directed cycle.  Their number is the necklace count
\begin{equation}
  N_q(k)=\frac1k\sum_{d\mid k}\varphi(d)q^{k/d}.
  \label{eq:necklace}
\end{equation}
Every decycling set meets every PCR.  The classical construction meets this
lower bound, so the ordinary minimum cardinality is $N_q(k)$
\cite{Mykkeltveit1972,ChamparnaudHanselPerrin2004}.  Consequently an MDS
selects exactly one distinct word from each PCR.  The converse is false:
one-per-PCR is a local cardinality condition and can leave cycles composed
of arcs from several PCRs.  The binary homopolymers $0^k$ and $1^k$ are
singleton PCR loops and are mandatory selections.

\subsection{Firing convention}

For a lower word $u\in\A^{k-1}$, write
\[
  \operatorname{lc}(u)=\{au:a\in\A\},
  \qquad
  \operatorname{rc}(u)=\{ua:a\in\A\}.
\]
Here ``lc'' and ``rc'' mean left and right companion sets; the latter is not
the reverse-complement operator $\RC$.  A forward firing, or $F$-move, is
valid at $u$ for a selected set $M$ precisely when
\[
  \operatorname{lc}(u)\subseteq M,
  \qquad
  F_u(M)=\bigl(M\setminus\operatorname{lc}(u)\bigr)
           \cup\operatorname{rc}(u).
\]
These are literal sets: any overlap between the companion sets, including
the homopolymer-loop case, is retained rather than counted twice.  The fact
that a valid $F$-move preserves ordinary-MDS status is the source-move result
of Mar\c{c}ais--DeBlasio--Kingsford
\cite[Section~4.1, Proposition~2]{MarcaisDeBlasioKingsford2024}.  We import
only this preservation fact, not any conjectural connectivity statement.

For an integer function $q$ on lower states, our gradient sign is
\begin{equation}
  (\nabla q)(w)=q(t(w))-q(h(w)).
  \label{eq:gradient}
\end{equation}
With this convention, a valid firing at $u$ changes the selected-edge
indicator by $\nabla\ind_u$; the set-overlap convention makes this equality
valid also at all homopolymer loops (the two loops in the binary
specialization).

\subsection{What is imported and what is proved here}

We import the ordinary cardinality theorem and valid-$F$-move preservation.
We prove the odd obstruction, arbitrary-zero-tie acyclicity of the chosen
cyclotomic seed, equivariant tie selection, the exact coboundary identity,
and the terminating defect descent.  The 2024 k-nonical paper supplies a
direct dated formulation of the symmetry-constrained problem
\cite{MarcaisElderKingsford2024}; it states that the minimum cardinality under
the symmetry constraint was not known and uses integer linear programming for
bounded orders.  Our result settles the binary ordinary-necklace-bound
existence question at every order.  It does not replace that paper's
optimization, implementation, or window-guarantee analysis, and its bounded
computations are not used as a theorem in our proof.
